\documentclass[12pt]{article}

% Настройка русского языка и шрифтов
\usepackage[utf8]{inputenc}
\usepackage[russian]{babel}
\usepackage[T2A]{fontenc}

% Математические пакеты
\usepackage{amsmath, amssymb, amsfonts, amsthm}

% Настройка страницы
\usepackage{geometry}
\geometry{a4paper, margin=0.8in}

% Другие полезные пакеты
\usepackage{hyperref}
\usepackage{booktabs}
\usepackage{enumitem}
\usepackage{longtable}
\usepackage{graphicx}
\usepackage{xcolor}
\usepackage{tcolorbox}

% Определение стилей для комментариев
\definecolor{commentgray}{rgb}{0.4,0.4,0.4}
\definecolor{antigravityblue}{rgb}{0.0, 0.45, 0.74}
\definecolor{theorygold}{rgb}{0.85, 0.65, 0.13}

\newcommand{\commentary}[1]{{\small\itshape\color{commentgray} #1}}
\newcommand{\antigravity}[1]{{\small\color{antigravityblue}\textbf{Antigravity:} #1}}

% Настройка ссылок
\hypersetup{
    colorlinks=true,
    linkcolor=blue,
    filecolor=magenta,      
    urlcolor=cyan,
    pdftitle={Мастер-разбор: Линейные модели},
}

\begin{document}

\section{Типология задач машинного обучения}
В обучении с учителем (\textit{Supervised Learning}) мы выделяем следующие основные задачи:

\begin{itemize}
    \item \textbf{Регрессия (Regression):} Целевая переменная $y \in \mathbb{R}$ является непрерывной. Мы ищем функциональную зависимость, которая минимизирует ошибку предсказания вещественного числа.
    \begin{itemize}
        \item \textit{Примеры:} Прогноз стоимости недвижимости на основе характеристик (площадь, район), предсказание температуры на завтра, оценка времени прибытия курьера (ETA).
        \item \textit{Метрики:} MSE, MAE, $R^2$.
    \end{itemize}

    \item \textbf{Бинарная классификация (Binary Classification):} $y \in \{0, 1\}$ (или $\{-1, 1\}$). Необходимо разделить объекты на две непересекающиеся группы. Модель обычно выдает вероятность $p \in [0, 1]$.
    \begin{itemize}
        \item \textit{Примеры:} Спам-фильтр (письмо спам или нет), кредитный скоринг (вернет ли клиент кредит), медицинская диагностика (наличие или отсутствие заболевания).
        \item \textit{Метрики:} Accuracy, Precision, Recall, F1-score, ROC-AUC.
    \end{itemize}

    \item \textbf{Многоклассовая классификация (Multiclass Classification):} $y \in \{1, \dots, K\}$. Выбор одного из $K$ взаимоисключающих классов.
    \begin{itemize}
        \item \textit{Примеры:} Распознавание рукописных цифр (0-9), классификация типов ирисов (Setosa, Versicolour, Virginica), определение языка текста.
        \item \textit{Подходы:} One-vs-Rest, One-vs-One, Softmax.
    \end{itemize}

    \item \textbf{Ранжирование (Ranking):} Задача не в предсказании точного значения, а в восстановлении правильного \textbf{порядка} объектов.
    \begin{itemize}
        \item \textit{Примеры:} Выдача поисковой системы (релевантные ссылки выше), ранжирование товаров в ленте рекомендаций, сортировка кандидатов в системе подбора персонала.
        \item \textit{Метрики:} nDCG, MRR, Precision@k.
    \end{itemize}
\end{itemize}

\section{Математический аппарат и обозначения}

\subsection{Обозначения}
Для формального описания линейных моделей введем следующие обозначения:
\begin{itemize}
    \item $\mathbf{X} \in \mathbb{R}^{N \times D}$ — матрица «объекты-признаки», где $N$ — число объектов, $D$ — число признаков.
    \item $\mathbf{y} \in \mathbb{R}^N$ — вектор истинных ответов (таргетов).
    \item $\mathbf{\omega} \in \mathbb{R}^D$ — вектор параметров (весов) модели.
    \item Предсказание модели: $\mathbf{\hat{y}} = \mathbf{X}\mathbf{\omega}$.
\end{itemize}

\subsection{Решение МНК (Ordinary Least Squares)}
Мы минимизируем сумму квадратов ошибок:
\[ Q(\mathbf{\omega}) = \|\mathbf{X}\mathbf{\omega} - \mathbf{y}\|^2_2 = (\mathbf{X}\mathbf{\omega} - \mathbf{y})^T (\mathbf{X}\mathbf{\omega} - \mathbf{y}) \]
Аналитическое решение находится через приравнивание градиента к нулю:
\[ \nabla_{\mathbf{\omega}} Q(\mathbf{\omega}) = 2\mathbf{X}^T \mathbf{X} \mathbf{\omega} - 2\mathbf{X}^T \mathbf{y} = 0 \]
Откуда получаем нормальное уравнение:
\[ \mathbf{\omega} = (\mathbf{X}^T \mathbf{X})^{-1} \mathbf{X}^T \mathbf{y} \]

\begin{tcolorbox}[colback=theorygold!5!white,colframe=theorygold!75!black,title=Условие существования]
    Решение существует и единственно, если матрица $\mathbf{X}^T \mathbf{X}$ невырождена (имеет полный ранг). Это эквивалентно линейной независимости признаков (отсутствию мультиколлинеарности).
\end{tcolorbox}

\subsection{Численные методы: QR-разложение}
Прямое вычисление $(\mathbf{X}^T \mathbf{X})^{-1}$ вычислительно дорого $O(D^3)$ и неустойчиво. Метод \textbf{QR-разложения} позволяет решить задачу стабильнее:
\begin{enumerate}
    \item Раскладываем $\mathbf{X} = \mathbf{QR}$, где $\mathbf{Q}$ — ортогональная матрица ($\mathbf{Q}^T\mathbf{Q} = \mathbf{I}$), а $\mathbf{R}$ — верхнетреугольная.
    \item Подставляем в нормальное уравнение: $(\mathbf{R}^T \mathbf{Q}^T \mathbf{Q} \mathbf{R}) \mathbf{\omega} = \mathbf{R}^T \mathbf{Q}^T \mathbf{y}$.
    \item Упрощаем: $\mathbf{R}^T \mathbf{R} \mathbf{\omega} = \mathbf{R}^T \mathbf{Q}^T \mathbf{y} \implies \mathbf{R} \mathbf{\omega} = \mathbf{Q}^T \mathbf{y}$.
    \item Система с треугольной матрицей $\mathbf{R}$ решается методом обратной подстановки за $O(D^2)$.
\end{enumerate}

\subsection{Градиентный спуск (Gradient Descent)}
Если аналитическое решение невозможно (огромные данные) или функция потерь неквадратична, используется итерационный подход:
\[ \mathbf{\omega}_{t+1} = \mathbf{\omega}_t - \eta \nabla Q(\mathbf{\omega}_t) \]
Где $\eta$ — шаг обучения (\textit{learning rate}).
Для OLS шаг градиентного спуска выглядит так:
\[ \mathbf{\omega}_{t+1} = \mathbf{\omega}_t - 2\eta \mathbf{X}^T (\mathbf{X}\mathbf{\omega}_t - \mathbf{y}) \]
Градиент направлен в сторону наискорейшего роста функции, поэтому мы вычитаем его, чтобы двигаться к минимуму.

\section{Обработка категориальных признаков}
Линейные модели требуют численного представления данных. Для работы с качественными переменными (город, профессия) применяется кодирование.

\subsection{One-Hot Encoding (OHE)}
Метод заменяет один признак с $M$ категориями на $M$ новых бинарных столбцов.
\begin{itemize}
    \item \textbf{Где можно использовать:}
    \begin{itemize}
        \item Категории не имеют естественного порядка (номинальные данные). 
        \item \textit{Пример:} Цвета автомобиля (красный, синий, зеленый), страны, типы почтовых отправлений.
        \item Линейные модели и нейросети — OHE позволяет модели выучить независимый вес для каждой категории.
    \end{itemize}
    \item \textbf{Где нельзя или не рекомендуется использовать:}
    \begin{itemize}
        \item \textbf{Высокая кардинальность:} Если у признака тысячи уникальных значений (напр., ID пользователя), OHE создаст разреженную матрицу огромного размера, что приведет к переобучению и нехватке памяти.
        \item \textbf{Деревья решений:} OHE "размывает" информацию. Дереву проще работать с порядковым кодированием (\textit{Label Encoding}), так как оно может делать сплиты по группам категорий.
        \item \textbf{Порядковые данные:} Если категории имеют смысл "больше-меньше" (напр., размер одежды S, M, L), лучше использовать \textit{Ordinal Encoding}, чтобы сохранить информацию о порядке.
    \end{itemize}
\end{itemize}

\begin{tcolorbox}[colback=red!5!white,colframe=red!75!black,title=Dummy Variable Trap]
    Для линейных моделей с интерсептом ($w_0$) необходимо удалять один столбец после OHE ($M-1$ признаков). В противном случае возникает строгая мультиколлинеарность, так как сумма всех OHE-столбцов всегда равна единице.
\end{tcolorbox}

\section{Интерпретируемость моделей}
Это способность модели объяснить, на основе каких признаков и почему было принято решение. По уровню прозрачности алгоритмы делятся на \textit{White-box} (прозрачные) и \textit{Black-box} (непрозрачные).

\subsection{Линейные модели}
Линейные модели относятся к наиболее интерпретируемым алгоритмам.
\begin{itemize}
    \item \textbf{Веса $\omega_i$:} Показывают вклад признака $x_i$. Если $\omega_i > 0$, рост признака увеличивает предсказание. Если $\omega_i < 0$ — уменьшает.
    \item \textbf{Абсолютное значение $|\omega_i|$:} Если признаки стандартизированы, то чем больше модуль веса, тем важнее признак для модели (\textit{Feature Importance}).
    \item \textit{Ограничение:} При наличии мультиколлинеарности веса теряют физический смысл и становятся неинтерпретируемыми.
\end{itemize}

\subsection{Решающие деревья}
Деревья интерпретируются через их структуру и правила принятия решений.
\begin{itemize}
    \item \textbf{Путь от корня к листу:} Представляет собой цепочку логических условий (IF-THEN), понятных человеку без специальной подготовки.
    \item \textbf{Feature Importance:} Рассчитывается на основе того, насколько сильно каждый признак уменьшает неопределенность (энтропию или критерий Джини) во всех узлах дерева.
    \item \textit{Ограничение:} Глубокие деревья («переобученные») становятся слишком сложными для визуального анализа человеком.
\end{itemize}

\section{Свойства и ограничения линейных моделей}
Понимание сильных и слабых сторон линейных алгоритмов необходимо для правильного выбора модели и качественной подготовки признаков.

\subsection{Свойства}
\begin{itemize}
    \item \textbf{Простота и скорость:} Обучение и предсказание выполняются крайне быстро, что критично для высоконагруженных систем.
    \item \textbf{Интерпретируемость:} Прямая связь между признаками и целевой переменной через веса.
    \item \textbf{Работа с редкими данными:} Хорошо справляются с разреженными матрицами (sparse data) после OHE.
\end{itemize}

\subsection{Ограничения}
\begin{itemize}
    \item \textbf{Линейность зависимости:} Модель не может выучить нелинейные паттерны и взаимодействия признаков (напр., $x_1 \cdot x_2$) без их явного добавления вручную (\textit{Polynomial Features}).
    \item \textbf{Чувствительность к выбросам:} Квадратичная ошибка в МНК придает аномалиям слишком большой вес, что может "развернуть" разделяющую плоскость. Решение — использование Huber Loss или MAE.
    \item \textbf{Мультиколлинеарность:} При наличии сильно коррелирующих признаков веса становятся огромными и неустойчивыми. Требуется регуляризация.
    \item \textbf{Масштаб признаков:} Веса напрямую зависят от единиц измерения. Для корректной работы регуляризации и интерпретации обязательна нормализация.
\end{itemize}

\section{Сведение обучения к задаче оптимизации}
Процесс «обучения» модели в ML формализуется как поиск параметров, минимизирующих некоторую функцию ошибки на обучающей выборке.

\subsection{Функция потерь и Функционал качества}
Важно различать ошибку на одном объекте и на всей выборке:
\begin{itemize}
    \item \textbf{Функция потерь $L(a(x), y)$:} Определяет величину ошибки на \textit{одном} конкретном объекте $x$. 
    \begin{itemize}
        \item Для регрессии: $(a(x) - y)^2$ (квадратичная), $|a(x) - y|$ (абсолютная).
        \item Для классификации: $\log(1 + \exp(-y \cdot a(x)))$ (логистическая).
    \end{itemize}
    \item \textbf{Функционал качества $Q(w)$:} Суммарная ошибка на всей обучающей выборке $X$:
    \[ Q(w) = \frac{1}{N} \sum_{i=1}^N L(a(x_i, w), y_i) \]
    Задача обучения — найти $w^* = \arg\min_w Q(w)$.
\end{itemize}

\subsection{Stochastic Gradient Descent (SGD) и Mini-batches}
Когда данных слишком много для вычисления градиента по всей выборке (Full Batch), используются стохастические методы:
\begin{enumerate}
    \item \textbf{SGD:} На каждом шаге выбирается один случайный объект $i$, и веса обновляются по его градиенту: $w \leftarrow w - \eta \nabla L(x_i)$. Это очень быстро, но траектория обучения сильно "шумит".
    \item \textbf{Mini-batch GD:} Золотая середина. Мы выбираем небольшой подмножество данных (\textit{batch}) размера $B$ (напр., 32, 64, 128) и вычисляем средний градиент по ним:
    \[ \nabla Q_{batch} = \frac{1}{B} \sum_{j=1}^B \nabla L(x_j) \]
    Это позволяет эффективно использовать векторизацию вычислений на CPU/GPU и уменьшить шум по сравнению с чистым SGD.
\end{enumerate}


\section{Регуляризация и борьба с мультиколлинеарностью}

\subsection{Мультиколлинеарность}
Это ситуация, когда признаки в матрице $\mathbf{X}$ сильно коррелируют друг с другом.
\begin{itemize}
    \item \textbf{Последствия:} Матрица $\mathbf{X}^T \mathbf{X}$ становится близкой к вырожденной. Это приводит к неустойчивости оценок весов: малейшее изменение в данных вызывает гигантские скачки значений $\omega_i$.
    \item \textbf{Решение:} Добавление штрафа за величину весов в функционал качества.
\end{itemize}

\subsection{L2-регуляризация (Ridge Regression)}
Добавляем квадрат $L2$-нормы весов: $Q_{reg}(\omega) = Q(\omega) + \lambda \|\omega\|^2_2$.
\begin{itemize}
    \item \textbf{Аналитическое решение:} $\omega = (\mathbf{X}^T \mathbf{X} + \lambda \mathbf{I})^{-1} \mathbf{X}^T \mathbf{y}$.
    \item \textbf{Плюсы:} Всегда делает матрицу обратимой, уменьшает дисперсию весов, имеет аналитическое решение.
    \item \textbf{Минусы:} Не зануляет веса (не делает отбор признаков).
\end{itemize}

\subsection{L1-регуляризация (Lasso)}
Добавляем $L1$-норму весов: $Q_{reg}(\omega) = Q(\omega) + \lambda \|\omega\|_1$.
\begin{itemize}
    \item \textbf{Особенность:} Из-за "острого" края ромба (линии уровня $L1$) точка оптимума часто попадает на оси координат.
    \item \textbf{Плюсы:} \textbf{Отбор признаков} (\textit{Feature Selection}) — зануляет веса у неважных признаков.
    \item \textbf{Минусы:} Нет аналитического решения, нестабильна при сильной корреляции (выбирает один случайный признак из группы коррелирующих).
\end{itemize}

\subsection{Сравнительная таблица}
\begin{table}[h!]
\centering
\begin{tabular}{@{}lll@{}}
\toprule
\textbf{Свойство} & \textbf{L2 (Ridge)} & \textbf{L1 (Lasso)} \\ \midrule
Штраф & $\sum \omega_i^2$ & $\sum |\omega_i|$ \\
Решение & Аналитическое & Численное \\
Отбор признаков & Нет (веса малы, но $\neq 0$) & Да (зануляет лишнее) \\
Устойчивость & Высокая & Низкая (чувствительна к шуму) \\
Интерпретация & Сохраняет все признаки & Дает разреженную модель \\ \bottomrule
\end{tabular}
\end{table}

\section{Логистическая регрессия}

\subsection{Сигмоида ($\sigma$)}
Для перевода выхода линейной модели $\omega^T x \in \mathbb{R}$ в вероятность $p \in [0, 1]$ используется логистическая функция (сигмоида):
\[ \sigma(z) = \frac{1}{1 + e^{-z}} \]
\begin{itemize}
    \item \textbf{Свойство:} $\sigma'(z) = \sigma(z)(1 - \sigma(z))$, что удобно при расчете градиентов.
    \item \textbf{Смысл:} Она "сплющивает" бесконечную прямую в интервал $(0, 1)$, позволяя интерпретировать результат как уверенность модели.
\end{itemize}

\subsection{Метод максимального правдоподобия (MLE)}
Мы хотим подобрать такие $\omega$, при которых наблюдаемые данные наиболее вероятны. Предположим, что $y_i \sim \text{Bernoulli}(p_i)$. Тогда правдоподобие всей выборки:
\[ L(\mathbf{\omega}) = \prod_{i=1}^n p_i^{y_i} (1-p_i)^{1-y_i} \]
Логарифмируя и меняя знак, переходим к минимизации \textbf{Log-Loss}:
\[ Q(\omega) = -\frac{1}{n} \sum_{i=1}^n [y_i \ln(\sigma(\omega^T x_i)) + (1-y_i) \ln(1 - \sigma(\omega^T x_i))] \]

\subsection{Изменения пространства признаков}
Линейные модели строят линейную разделяющую поверхность в текущем пространстве признаков. Чтобы выучить нелинейные зависимости, мы можем трансформировать исходное пространство:
\begin{itemize}
    \item \textbf{Полиномиальные признаки:} Добавление степеней и произведений признаков $(x_1, x_2 \to x_1^2, x_1x_2, x_2^2)$. Это позволяет модели строить криволинейные границы (эллипсы, параболы).
    \item \textbf{RBF и Ядра:} Переход в бесконечномерное пространство, где данные становятся линейно разделимыми (основа SVM с ядрами).
    \item \textbf{Feature Engineering:} Ручное создание признаков (напр., $\sin(x)$, $\text{sign}(x)$) на основе знаний о предметной области.
\end{itemize}

\section{Функции потерь классификации}
Для задач классификации с метками $y \in \{-1, 1\}$ часто используют функции от \textit{отступа} (margin) $M = y \cdot (\omega^T x)$.

\subsection{Perceptron Loss}
\[ L(M) = \max(0, -M) \]
\begin{itemize}
    \item Ошибкой считаются только те объекты, которые лежат на «чужой» стороне разделяющей плоскости.
    \item Не дает никакой уверенности в классификации — объект может быть очень близко к границе.
\end{itemize}

\subsection{Hinge Loss и SVM}
\[ L(M) = \max(0, 1 - M) \]
Это стандартная функция потерь для \textbf{SVM (Метод опорных векторов)}.
\begin{itemize}
    \item \textbf{Зазор (Margin):} В отличие от перцептрона, Hinge Loss требует, чтобы объекты не просто были классифицированы верно ($M>0$), но и находились на расстоянии хотя бы 1 от разделяющей прямой ($M \ge 1$).
    \item Все объекты с $M < 1$ вносят вклад в ошибку. Это обеспечивает «максимальный разделяющий зазор», что делает модель более устойчивой к шуму.
\end{itemize}

\newpage
\section{Kernel Trick и ядерные методы}
Ядерный переход (\textit{Kernel Trick}) позволяет линейным моделям работать в нелинейных пространствах без явного вычисления координат новых признаков.

\subsection{Суть метода}
Многие алгоритмы (SVM, PCA, Линейная регрессия) зависят только от скалярных произведений объектов $\langle x_i, x_j \rangle$. Мы можем заменить их на \textbf{ядерную функцию} $K(x_i, x_j) = \langle \phi(x_i), \phi(x_j) \rangle$, которая вычисляет скалярное произведение в некотором пространстве более высокой размерности $\phi(\cdot)$.

\subsection{Виды ядер}
\begin{itemize}
    \item \textbf{Линейное (Linear):} $K(x, y) = x^T y$. Обычное скалярное произведение. Используется, когда данных много или они уже линейно разделимы.
    \item \textbf{Полиномиальное (Polynomial):} $K(x, y) = (x^T y + c)^d$. Позволяет учитывать взаимодействия признаков до степени $d$.
    \item \textbf{RBF / Гауссово (Radial Basis Function):} $K(x, y) = \exp(-\gamma \|x - y\|^2)$. Бесконечномерное ядро. Самое популярное для нелинейных задач. $\gamma$ определяет "радиус влияния" каждого объекта.
    \item \textbf{Сигмоидальное (Sigmoid):} $K(x, y) = \tanh(\alpha x^T y + c)$. Позволяет имитировать поведение двухслойной нейронной сети.
\end{itemize}

\subsection{Kernel Density Estimation (KDE)}
Ядерная оценка плотности — это непараметрический способ оценки функции распределения случайной величины.
\[ \hat{f}_h(x) = \frac{1}{nh} \sum_{i=1}^n K\left(\frac{x - x_i}{h}\right) \]
\begin{itemize}
    \item Вместо грубых "ящиков" гистограммы мы ставим "колокол" ядра (обычно Гауссова) над каждой точкой данных.
    \item Параметр \textit{bandwidth} ($h$) контролирует гладкость: слишком малое $h$ дает шум, слишком большое — переглаженное распределение.
\end{itemize}

\newpage
\section{Многоклассовая классификация}
Когда число классов $K > 2$, бинарная логистическая регрессия обобщается до многоклассовой (Multinomial Logistic Regression).

\subsection{Вероятности и Softmax}
Вместо одного числа модель выдает вектор "сырых" ответов (\textit{logits}) $z = (z_1, \dots, z_K)$. Чтобы превратить их в вероятности, используется функция \textbf{Softmax}:
\[ P(y=k|x) = \frac{e^{z_k}}{\sum_{j=1}^K e^{z_j}} \]
\begin{itemize}
    \item Сумма всех вероятностей всегда равна 1.
    \item Каждый компонент $P \in (0, 1)$.
    \item Softmax является многомерным обобщением сигмоиды.
\end{itemize}

\subsection{Функция потерь: Cross-Entropy}
Для многоклассовой классификации минимизируется кросс-энтропия между истинным распределением (One-Hot вектор $y$) и предсказанным ($p$):
\[ Q(\omega) = -\sum_{i=1}^n \sum_{k=1}^K y_{ik} \ln p_{ik} \]
Если объект относится к классу $k$, то $y_{ik}=1$, и мы просто минимизируем $-\ln p_{ik}$ (максимизируем вероятность правильного класса).

\subsection{Алгоритмы сведения к бинарным задачам}
Если мы не хотим использовать Softmax напрямую, можно применить стратегии:
\begin{itemize}
    \item \textbf{One-vs-Rest (OvR):} Обучаем $K$ классификаторов. Каждый учится отличать свой класс от всех остальных вместе взятых. Предсказание — класс с максимальной уверенностью.
    \item \textbf{One-vs-One (OvO):} Обучаем $K(K-1)/2$ классификаторов для каждой пары классов. Итоговое решение принимается голосованием большинства. Более устойчива к дисбалансу, но требует много моделей.
\end{itemize}

\newpage
\section{Deep Q\&A: Вопросы и подробные ответы}

\begin{enumerate}[label=\textbf{Q\arabic*:}, leftmargin=1.5cm]
    \item \textbf{Почему в Ridge-регрессии веса никогда не становятся ровно нулями?} \\
    \textit{Ответ:} Линии уровня L2-штрафа — это сферы. Касание линий уровня функции потерь и сферы может произойти в любой точке. Вероятность того, что точка касания попадет точно на координатную ось в высокомерном пространстве, практически нулевая. В отличие от L1 (ромб), у L2 нет "углов" на осях.

    \item \textbf{Как мультиколлинеарность влияет на дисперсию весов?} \\
    \textit{Ответ:} Дисперсия весов обратно пропорциональна определителю $\mathbf{X}^T \mathbf{X}$. При мультиколлинеарности матрица близка к вырожденной, её определитель мал, что раздувает дисперсию оценок весов. Это делает модель крайне чувствительной к малейшим изменениям в обучающей выборке.

    \item \textbf{В чем разница между моделями One-vs-Rest и One-vs-One?} \\
    \textit{Ответ:} 
    \begin{itemize}
        \item \textbf{OvR:} Обучаем $K$ классификаторов ("Класс $i$ против всех остальных"). Быстрее, но может возникнуть дисбаланс классов.
        \item \textbf{OvO:} Обучаем $K(K-1)/2$ классификаторов для каждой пары. Более устойчива к дисбалансу, но вычислительно дороже при большом числе классов.
    \end{itemize}

    \item \textbf{Зачем нужна калибровка вероятностей (напр. по Платту)?} \\
    \textit{Ответ:} Некоторые модели (например, SVM) не выдают вероятности напрямую. Сигмоидальная калибровка Платта обучает логистическую регрессию над выходами модели, чтобы превратить их в честные доверительные интервалы $[0, 1]$.

    \item \textbf{Что такое Rank Deficiency и как его лечить?} \\
    \textit{Ответ:} Это ситуация, когда число признаков $D$ больше числа объектов $N$, либо признаки линейно зависимы. В этом случае $\mathbf{X}^T \mathbf{X}$ необратима. Лечится добавлением L2-регуляризации, которая делает матрицу $(\mathbf{X}^T \mathbf{X} + \lambda I)$ строго положительно определенной и всегда обратимой.
\end{enumerate}

\vspace{1cm}
\antigravity{Этот уровень проработки теории отличает Senior ML-инженера от новичка. Помните: любая сложная нейросеть — это всего лишь композиция простых линейных преобразований и нелинейностей.}

\end{document}
